% Tutoriel cellyod — compiler avec : latexmk -lualatex tutorial_cellyod.tex
% (figures et données : ./make_figures.sh)
\documentclass[10pt,a4paper]{article}

\usepackage{fontspec}
\usepackage[french]{babel}
\usepackage{amsmath,amssymb}
\usepackage{booktabs,tabularx}
\usepackage{graphicx}
\usepackage{subcaption}
\usepackage{float}
\usepackage{tikz}
\usetikzlibrary{arrows.meta,calc,positioning,decorations.pathreplacing}
\usepackage{pgfplots}
\pgfplotsset{compat=1.18}
\usepackage[margin=2.1cm]{geometry}
\usepackage{xcolor}
\usepackage{listings}
\usepackage{enumitem}
\usepackage{multicol}
\usepackage[colorlinks=true,linkcolor=blue!50!black,urlcolor=blue!50!black]{hyperref}

\setlist{itemsep=1pt,topsep=3pt}
\emergencystretch=3em
\definecolor{cell}{HTML}{0F111A}
\definecolor{bundle}{HTML}{F07470}
\definecolor{accent}{HTML}{2D7DD2}
\definecolor{seed}{HTML}{8E8F9C}

\lstset{basicstyle=\ttfamily\footnotesize,columns=fullflexible,frame=single,framerule=0.3pt,
  backgroundcolor=\color{gray!6},breaklines=true,keepspaces=true,upquote=true}

\newcommand{\code}[1]{\texttt{#1}}
\newcommand{\cy}{\code{cellyod}}
\newcommand{\cyg}{\code{cellyod-gui}}
\newcommand{\vect}[1]{\boldsymbol{#1}}

\title{\cy{}\\[3pt]
  \large Génération de tissus cellulaires par Voronoï et itérations de Lloyd :\\
  tutoriel, méthode et organisation du code}
\author{Dépôt \code{cellYOd}}
\date{\today}

\begin{document}
\maketitle

\begin{abstract}
  \cy{} pave un domaine plan par un diagramme de Voronoï que des itérations de Lloyd régularisent. Une fraction
  des cellules sont des \emph{bundles} (faisceaux vasculaires) qui attirent les autres au sens d'une distance de
  Mahalanobis : selon deux paramètres de forme, l'intensité $\alpha$ de l'attraction et l'anisotropie $K$ de la
  métrique, le tissu passe continûment du pavage hexagonal régulier à des tissus très hétérogènes, organisés en
  amas ou en colonnes. Le domaine peut avoir une forme quelconque, définie par une fonction $f(x,y)$. Le
  programme comprend un cœur C++17 sans dépendance, une ligne de commande (\cy) et une interface graphique
  (\cyg, raylib) où tous les paramètres se règlent à la souris pendant que le tissu évolue. Ce document montre
  comment s'en servir (section~\ref{sec:prise}), décrit la méthode en détail (section~\ref{sec:methode}),
  les fichiers (section~\ref{sec:fichiers}), l'interface (section~\ref{sec:gui}), des études types
  (section~\ref{sec:etudes}) et l'organisation du code (section~\ref{sec:code}).
\end{abstract}

\tableofcontents

%======================================================================
\section{Prise en main}\label{sec:prise}
%======================================================================

\subsection{Compilation}

\begin{lstlisting}
make                 # cellyod (ligne de commande) et cellyod-gui (interface)
make cellyod         # ligne de commande seule : aucune dépendance, simple compilateur C++17
\end{lstlisting}
L'interface demande raylib (\code{brew install raylib}) ; raygui, rini et la police d'interface sont fournis
dans \code{gui/}. Le cœur et la ligne de commande n'utilisent que la bibliothèque standard.

\subsection{Trois façons de l'utiliser}

\paragraph{Explorer avec l'interface.} C'est le point de départ conseillé : on voit le tissu se former
itération par itération et on observe immédiatement l'effet de chaque paramètre.
\begin{lstlisting}
./cellyod-gui                        # paramètres par défaut (réseau hexagonal perturbé, Lloyd seul)
./cellyod-gui examples/tissu.txt     # fichier de paramètres... ou fichier de cellules déjà calculé
\end{lstlisting}
Dans le panneau de droite, mettre par exemple \code{bundleFraction} à 0.05 (les germes sont régénérés),
puis faire varier \code{alpha} et \code{anisotropy} pendant que les itérations tournent (bouton \emph{Lancer}
ou barre d'espace). Le bouton \emph{Enregistrer} écrit le fichier de cellules et sa vue SVG ;
\emph{Paramètres} écrit un fichier de paramètres relisible par la ligne de commande.

\paragraph{Produire un tissu en ligne de commande.} Un fichier de paramètres « mot-clé valeur » décrit le
tissu ; \code{mot-clé=valeur} sur la ligne de commande remplace une valeur du fichier :
\begin{lstlisting}
cd examples
../cellyod tissu.txt                                  # -> tissu_cells.txt et tissu.svg
../cellyod tissu.txt alpha=0.1 "anisotropy=1000 90"  # variante, sans modifier le fichier
../cellyod tissu.txt svg=none output=a.txt            # sans SVG
../cellyod alpha=0.05 bundleFraction=0.05 Lx=60 Ly=30 # sans fichier : défauts + remplacements
\end{lstlisting}
La sortie résume le tissu obtenu :
\lstinputlisting{figures/sortie_tissu.txt}

\paragraph{Reprendre un calcul.} Un fichier de cellules contient les paramètres et les germes : il suffit pour
poursuivre les itérations, en changeant éventuellement les paramètres « vivants » (section~\ref{sec:effets}) :
\begin{lstlisting}
../cellyod -load tissu_cells.txt lloydIterations=200 output=tissu200.txt
../cellyod -load tissu_cells.txt lloydIterations=80 alpha=0.1   # 30 itérations de plus avec alpha = 0.1
\end{lstlisting}
Une reprise donne exactement les mêmes germes qu'un calcul d'une traite (section~\ref{sec:reprise}).

\subsection{Exemples fournis}

\begin{center}
\begin{tabularx}{\linewidth}{@{}lXl@{}}
  \toprule
  Fichier & Tissu & Réglage principal \\
  \midrule
  \code{examples/tissu.txt} & référence de l'article : amas autour des bundles &
    \code{random}, 5\,\% bundles, $\alpha=0.05$, $K=1$ \\
  \code{examples/colonnes.txt} & organisation en colonnes & idem, $K=1000$ \\
  \code{examples/tomographie.txt} & coupe tomographique & \code{square}, \code{disorder 0}, 5 it., $\alpha=0.02$,
    $K=10^4$ \\
  \code{examples/rein.txt} & domaine non convexe (disque échancré) & \code{shape}, référence \\
  \bottomrule
\end{tabularx}
\end{center}

\subsection{Unités}

Toutes les grandeurs sont \textbf{sans dimension}. L'unité de longueur est le diamètre équivalent visé des
cellules : une cellule a une aire visée $A_0=\pi/4$, et l'espacement visé entre germes vaut
$s=\sqrt{A_0}=\sqrt{\pi}/2\approx 0.886$. Un rectangle $L_x\times L_y$ contient donc environ
$L_xL_y/A_0$ germes (2037 pour $40\times40$). Pour une application dimensionnée, il suffit de multiplier toutes les
coordonnées par le diamètre de cellule voulu.

%======================================================================
\section{Paramètres}\label{sec:params}
%======================================================================

Un fichier de paramètres contient une ligne par mot-clé : \code{mot-clé valeur(s)}, \code{\#} commençant un
commentaire. Les mots-clés absents gardent leur valeur par défaut ; un mot-clé inconnu ou une valeur invalide
sont refusés avec un message qui donne le fichier et la ligne. La liste complète est donnée par
\code{cellyod -h} :
\lstinputlisting{figures/aide.txt}

\subsection{Effet d'une modification}\label{sec:effets}

Chaque paramètre est déclaré soit \emph{de régénération} (marqué \code{*} dans l'interface), soit
\emph{vivant} :
\begin{itemize}
  \item \textbf{régénération} : \code{Lx}, \code{Ly}, \code{shape}, \code{lattice}, \code{disorder},
    \code{seed}, \code{bundleFraction}. Ils définissent l'état initial (germes et choix des bundles) ; les
    modifier recommence le calcul à l'itération 0 ;
  \item \textbf{vivants} : \code{beta}, \code{lloydIterations}, \code{alpha}, \code{anisotropy},
    \code{anisotropyAngle}, \code{bundleMotion}. Ils n'interviennent que dans la règle d'itération : une
    modification s'applique à l'itération suivante, sans repartir de zéro. On peut donc enchaîner plusieurs
    régimes (par exemple quelques itérations de Lloyd pur, puis l'attraction), dans l'interface ou avec
    \code{-load}.
\end{itemize}

%======================================================================
\section{Méthode}\label{sec:methode}
%======================================================================

Le calcul enchaîne : (1) placement des germes initiaux dans le rectangle $[0,L_x]\times[0,L_y]$ ; (2) choix des
bundles ; (3) itérations, chacune consistant à calculer le pavage de Voronoï des germes, à le restreindre à la
forme du domaine, puis à déplacer chaque germe ; (4) construction du pavage à sommets partagés et calcul des
métriques. Seul le pavage de Voronoï exact des germes est produit : il n'y a aucune retouche de la topologie
(fusion d'arêtes courtes, transitions T1\dots) comme dans \code{nodeGen}.

\subsection{Germes initiaux}\label{sec:germes}

Quatre états initiaux (\code{lattice}), tous à la densité visée (figure~\ref{fig:lattices}) :
\begin{itemize}
  \item \code{hex} : réseau hexagonal de pas $a=\sqrt{2A_0/\sqrt3}$, rangées espacées de $a\sqrt3/2$ et
    décalées d'un demi-pas une sur deux. Chaque germe est déplacé d'un vecteur de module
    $r=\text{\code{disorder}}\cdot a\sqrt{u_1}$ et d'angle $2\pi u_2$ ($u_1,u_2$ uniformes sur $[0,1]$), soit un
    déplacement uniforme dans le disque de rayon $\text{\code{disorder}}\cdot a$. Avec \code{disorder 0}, le pavage
    est un pavage hexagonal parfait ;
  \item \code{square} : même construction sur une grille carrée de pas $\sqrt{A_0}$ ; avec \code{disorder 0}, la
    grille carrée exacte utilisée dans l'article pour la comparaison aux coupes tomographiques ;
  \item \code{random} (alias \code{poisson}) : $N=\mathrm{round}(L_xL_y/A_0)$ points tirés uniformément avec une
    distance minimale $0.7\,s$ (tirage avec rejet, au plus $200N$ essais, puis complément sans contrainte) ;
  \item \code{uniform} : $N$ points tirés uniformément, sans distance minimale (processus de Poisson). Des germes
    très proches donnent des cellules minuscules : c'est l'état le plus désordonné.
\end{itemize}
Le générateur est un \code{mt19937\_64} initialisé par \code{seed} : à paramètres égaux, le résultat est
identique d'une exécution à l'autre ; pour les mêmes réglages, il donne les mêmes cellules et les mêmes
métriques que \code{nodeGenPlantCom}.

\begin{figure}[H]
  \centering
  \begin{subfigure}{0.24\linewidth}\includegraphics[width=\linewidth]{figures/lattice_hex.png}\caption{\code{hex}}\end{subfigure}\hfill
  \begin{subfigure}{0.24\linewidth}\includegraphics[width=\linewidth]{figures/lattice_square.png}\caption{\code{square}}\end{subfigure}\hfill
  \begin{subfigure}{0.24\linewidth}\includegraphics[width=\linewidth]{figures/lattice_random.png}\caption{\code{random}}\end{subfigure}\hfill
  \begin{subfigure}{0.24\linewidth}\includegraphics[width=\linewidth]{figures/lattice_uniform.png}\caption{\code{uniform}}\end{subfigure}
  \caption{Pavages de Voronoï des germes initiaux ($14\times14$, \code{disorder 0.3}, 0 itération).}
  \label{fig:lattices}
\end{figure}

\subsection{Bundles}

Une fraction \code{bundleFraction} des germes (parmi ceux dont la cellule est dans la forme) est tirée au hasard
comme bundles, avec le même générateur, après le placement des germes. S'y ajoute, pour chaque ligne
\code{bundle x y}, le germe le plus proche du point $(x,y)$. Dans l'interface, un clic sur une cellule en fait un
bundle (ou l'inverse) ; les bundles ajoutés à la souris sont mémorisés comme lignes \code{bundle x y}, de sorte
qu'ils survivent à une régénération et sont enregistrés avec les paramètres.

\subsection{Pavage de Voronoï}\label{sec:voronoi}

La cellule du germe $\vect p_i$ est l'ensemble des points plus proches de $\vect p_i$ que de tout autre germe :
c'est l'intersection des demi-plans $\{\vect x : (\vect x-\vect m_{ij})\cdot(\vect p_j-\vect p_i)\le 0\}$,
$\vect m_{ij}=(\vect p_i+\vect p_j)/2$. \cy{} la construit directement sous cette forme :
\begin{enumerate}
  \item les germes sont rangés dans une grille de cases de côté $s$ ;
  \item la cellule part du rectangle $[0,L_x]\times[0,L_y]$ et est découpée (algorithme de Sutherland--Hodgman
    pour un demi-plan) par les germes des couronnes de cases $k=0,1,2,\dots$ autour de $\vect p_i$ ;
  \item on s'arrête dès que $k\,s\ge 2R$, $R$ étant la distance de $\vect p_i$ au sommet le plus éloigné de la
    cellule courante : un germe plus lointain que $2R$ a sa médiatrice hors de la cellule et ne peut plus la
    couper.
\end{enumerate}
Le coût est proportionnel au nombre de germes (une vingtaine de voisins examinés par cellule), soit environ
4~ms par itération pour 2000 cellules. Les cellules obtenues sont convexes et orientées dans le sens
trigonométrique. Dans chaque case, les germes sont traités dans l'ordre de leurs coordonnées : le résultat,
arrondis compris, ne dépend pas de leur numérotation, ce qui garantit l'exactitude des reprises.

\subsection{Itérations de Lloyd}\label{sec:lloyd}

L'algorithme de Lloyd minimise l'énergie de la tessellation centroïdale de Voronoï
\begin{equation}
  E_{\mathrm{CVL}}(P)=\sum_i\int_{V_i}\|\vect x-\vect p_i\|^2\,d\vect x ,
\end{equation}
qui mesure la compacité des cellules : son minimum est atteint quand chaque germe est le centroïde
$\vect c_i$ de sa cellule $V_i$ (pour des cellules de même aire, les hexagones réguliers). Chaque itération
déplace simultanément tous les germes en $\vect p_i'=\vect c_i$ ; l'énergie décroît à chaque pas. Partant d'un
tirage uniforme, quelques itérations suffisent à rendre le pavage régulier (figure~\ref{fig:lloyd}) ; le
déplacement moyen \code{lastMove} tend alors vers 0.

\begin{figure}[H]
  \centering
  \begin{subfigure}{0.3\linewidth}\includegraphics[width=\linewidth]{figures/lloyd_0.png}\caption{0 itération}\end{subfigure}\hfill
  \begin{subfigure}{0.3\linewidth}\includegraphics[width=\linewidth]{figures/lloyd_5.png}\caption{5 itérations}\end{subfigure}\hfill
  \begin{subfigure}{0.3\linewidth}\includegraphics[width=\linewidth]{figures/lloyd_50.png}\caption{50 itérations}\end{subfigure}
  \caption{Régularisation de Lloyd seule ($\alpha=0$) à partir d'un tirage \code{uniform}.}
  \label{fig:lloyd}
\end{figure}

\subsection{Attraction vers les bundles}\label{sec:attraction}

Le modèle de Trivino \textit{et al.} (morphogenèse de Voronoï variationnelle) voit un tissu végétal comme deux
populations : les cellules ordinaires, qui remplissent l'espace, et les bundles, qui apportent eau et
nutriments. Le tissu minimise
\begin{equation}
  E(P,B)=E_{\mathrm{CVL}}(P)-\alpha\,E_{\mathrm{att}}(P,B),\qquad
  E_{\mathrm{att}}=\sum_i\|\vect p_i-\vect b_{j(i)}\|^2_M ,
\end{equation}
où $\vect b_{j(i)}$ est le bundle le plus proche de la cellule $i$ au sens de la distance de Mahalanobis
\begin{equation}
  d_M(\vect u,\vect v)=\sqrt{(\vect u-\vect v)^{\mathsf T}\Sigma^{-1}(\vect u-\vect v)} .
\end{equation}
\textbf{Règle d'itération.} Pour chaque germe ordinaire : (1) centroïde $\vect c_i$ de sa cellule ; (2) bundle
$\vect b$ le plus proche de $\vect c_i$ au sens de $d_M$ ; (3) nouveau germe
\begin{equation}\label{eq:update}
  \vect p_i'=\frac{\vect c_i+\alpha\,\vect b}{1+\alpha} ,
\end{equation}
c'est-à-dire $\vect c_i$ déplacé vers $\vect b$ d'une fraction $\alpha/(1+\alpha)$ du segment
$[\vect c_i,\vect b]$. C'est l'itération de point fixe de la fonctionnelle ; avec $\alpha=0$ on retrouve
exactement Lloyd. Tous les germes sont mis à jour simultanément, avec les positions des bundles du début de
l'itération. Les bundles eux-mêmes suivent \code{bundleMotion} : \code{lloyd} (défaut : vers le centroïde de leur
cellule, $B$ est optimisé avec $P$), \code{fixed} (immobiles) ou \code{attract} (attirés comme les autres
cellules, vers le bundle le plus proche autre qu'eux-mêmes).

\textbf{Métrique.} $\Sigma=K\,\vect a\vect a^{\mathsf T}+\vect b\vect b^{\mathsf T}$, avec $\vect a$ le vecteur
unitaire d'angle \code{anisotropyAngle} (direction des colonnes) et $\vect b$ son orthogonal ; avec l'angle par
défaut de $90^\circ$, $\Sigma=\mathrm{diag}(1,K)$ et $d_M^2=\Delta x^2+\Delta y^2/K$. Toute matrice définie positive
s'écrit ainsi à un facteur près, et seule la forme de $\Sigma$ compte (un facteur multiplie toutes les distances
sans changer le bundle le plus proche).

\textbf{Les deux paramètres de forme} (figure~\ref{fig:shapespace}) :
\begin{itemize}
  \item $\alpha$ règle surtout l'\emph{hétérogénéité des aires} : les cellules proches d'un bundle rétrécissent.
    Ordres de grandeur de l'article : $\alpha$ entre 0 et 0.2, référence 0.05 ;
  \item $K$ règle l'\emph{organisation spatiale} : $K=1$ (distance euclidienne) donne des amas de petites
    cellules autour de chaque bundle ; pour $K\gg1$, la distance est presque insensible aux écarts le long de
    $\vect a$ : un bundle attire toute la bande de cellules de même abscisse (pour l'angle $90^\circ$) et le tissu
    s'organise en bandes ($K=1000$ dans l'article, $10^4$ pour les coupes tomographiques).
\end{itemize}

\begin{figure}[H]
  \centering
  \setlength{\tabcolsep}{2pt}
  \begin{tabular}{c ccc}
    & $K=1$ & $K=100$ & $K=1000$\\
    \raisebox{0.13\linewidth}{$\alpha=0$} &
    \includegraphics[width=0.27\linewidth]{figures/shape_a0_K1.png} &
    \includegraphics[width=0.27\linewidth]{figures/shape_a0_K100.png} &
    \includegraphics[width=0.27\linewidth]{figures/shape_a0_K1000.png}\\
    \raisebox{0.13\linewidth}{$\alpha=0.05$} &
    \includegraphics[width=0.27\linewidth]{figures/shape_a0.05_K1.png} &
    \includegraphics[width=0.27\linewidth]{figures/shape_a0.05_K100.png} &
    \includegraphics[width=0.27\linewidth]{figures/shape_a0.05_K1000.png}\\
    \raisebox{0.13\linewidth}{$\alpha=0.15$} &
    \includegraphics[width=0.27\linewidth]{figures/shape_a0.15_K1.png} &
    \includegraphics[width=0.27\linewidth]{figures/shape_a0.15_K100.png} &
    \includegraphics[width=0.27\linewidth]{figures/shape_a0.15_K1000.png}
  \end{tabular}
  \caption{Espace des formes : $30\times30$, \code{random}, 5\,\% de bundles (points rouges), 50 itérations,
    même graine. Avec $\alpha=0$, $K$ est sans effet.}
  \label{fig:shapespace}
\end{figure}

\textbf{Pas de convergence.} À $\alpha$ et $K$ fixés, l'itération ne converge pas vers un état stationnaire :
les cellules continuent de se resserrer autour des bundles. Le nombre d'itérations est donc un
\emph{troisième paramètre de forme}, et deux tissus ne se comparent qu'à nombre d'itérations égal
(figures~\ref{fig:iter} et \ref{fig:curves}). Pour $K=1$, $\sigma_A/\langle A\rangle$ croît sans saturer
(0.30 après 5 itérations, 0.63 après 50, 0.75 après 200) tandis que le déplacement moyen décroît lentement ;
pour $K=1000$, elle passe par un maximum (environ 1.0 vers 50--80 itérations) puis décroît, et les germes ne
cessent pas de se déplacer (\code{lastMove} $\approx 0.3$ espacement par itération).

\begin{figure}[H]
  \centering
  \begin{subfigure}{0.24\linewidth}\includegraphics[width=\linewidth]{figures/iter_0.png}\caption{0}\end{subfigure}\hfill
  \begin{subfigure}{0.24\linewidth}\includegraphics[width=\linewidth]{figures/iter_10.png}\caption{10}\end{subfigure}\hfill
  \begin{subfigure}{0.24\linewidth}\includegraphics[width=\linewidth]{figures/iter_50.png}\caption{50}\end{subfigure}\hfill
  \begin{subfigure}{0.24\linewidth}\includegraphics[width=\linewidth]{figures/iter_200.png}\caption{200 itérations}\end{subfigure}
  \caption{Évolution du tissu de référence ($\alpha=0.05$, $K=1$) avec le nombre d'itérations.}
  \label{fig:iter}
\end{figure}

\begin{figure}[H]
  \centering
  \begin{tikzpicture}
    \begin{axis}[width=0.48\linewidth,height=5cm,xlabel={itérations},ylabel={$\sigma_A/\langle A\rangle$},
        legend pos=south east,legend style={font=\footnotesize},grid=major,grid style={gray!20}]
      \addplot[accent,thick,mark=*,mark size=1.2pt] table[x=iteration,y=cvA]{figures/iter_K1.dat};
      \addplot[bundle,thick,mark=square*,mark size=1.2pt] table[x=iteration,y=cvA]{figures/iter_K1000.dat};
      \legend{$K=1$,$K=1000$}
    \end{axis}
  \end{tikzpicture}\hfill
  \begin{tikzpicture}
    \begin{axis}[width=0.48\linewidth,height=5cm,xlabel={$\alpha$ (50 itérations)},ylabel={$\sigma_A/\langle A\rangle$},
        legend pos=south east,legend style={font=\footnotesize},grid=major,grid style={gray!20}]
      \addplot[accent,thick,mark=*,mark size=1.2pt] table[x=alpha,y=cvA]{figures/alpha_K1.dat};
      \addplot[bundle,thick,mark=square*,mark size=1.2pt] table[x=alpha,y=cvA]{figures/alpha_K1000.dat};
      \legend{$K=1$,$K=1000$}
    \end{axis}
  \end{tikzpicture}\\[2mm]
  \begin{tikzpicture}
    \begin{axis}[width=0.48\linewidth,height=5cm,xlabel={itérations},ylabel={$\langle SI\rangle$},
        legend pos=south east,legend style={font=\footnotesize},grid=major,grid style={gray!20}]
      \addplot[accent,thick,mark=*,mark size=1.2pt] table[x=iteration,y=meanSI]{figures/iter_K1.dat};
      \addplot[bundle,thick,mark=square*,mark size=1.2pt] table[x=iteration,y=meanSI]{figures/iter_K1000.dat};
      \legend{$K=1$,$K=1000$}
    \end{axis}
  \end{tikzpicture}\hfill
  \begin{tikzpicture}
    \begin{axis}[width=0.48\linewidth,height=5cm,xlabel={itérations},ylabel={\code{lastMove}},ymode=log,
        legend pos=south west,legend style={font=\footnotesize},grid=major,grid style={gray!20}]
      \addplot[accent,thick,mark=*,mark size=1.2pt] table[x=iteration,y=lastMove]{figures/iter_K1.dat};
      \addplot[bundle,thick,mark=square*,mark size=1.2pt] table[x=iteration,y=lastMove]{figures/iter_K1000.dat};
      \legend{$K=1$,$K=1000$}
    \end{axis}
  \end{tikzpicture}
  \caption{Métriques ($30\times30$, \code{random}, 5\,\% de bundles, $\alpha=0.05$ sauf en haut à droite) ;
    données produites par le fichier \code{morphology} (section~\ref{sec:morpho}).}
  \label{fig:curves}
\end{figure}

\begin{figure}[H]
  \centering
  \begin{subfigure}{0.3\linewidth}\includegraphics[width=\linewidth]{figures/angle30.png}\caption{$K=1000$, angle $30^\circ$}\end{subfigure}\hfill
  \begin{subfigure}{0.3\linewidth}\includegraphics[width=\linewidth]{figures/motion_fixed.png}\caption{\code{bundleMotion fixed}}\end{subfigure}\hfill
  \begin{subfigure}{0.3\linewidth}\includegraphics[width=\linewidth]{figures/tomo.png}\caption{coupe tomographique}\end{subfigure}
  \caption{Variantes : orientation des colonnes, bundles immobiles ($\alpha=0.05$, $K=1$, 50 itérations),
    réglage « coupe tomographique » (\code{square}, \code{disorder 0}, $\alpha=0.02$, $K=10^4$, 5 itérations).}
\end{figure}

\subsection{Forme du domaine}\label{sec:forme}

\textbf{Principe.} Les germes remplissent toujours le rectangle ; le tissu est la partie
$D=\{f(x,y)\le0\}$ de ce rectangle, $f$ étant donnée par l'expression \code{shape}. Le pavage de Voronoï est
calculé dans le rectangle, puis chaque cellule $V_i$ est restreinte à $V_i\cap D$ :
\begin{itemize}
  \item une cellule \emph{intérieure} ($V_i\subset D$) est inchangée et suit la règle~\eqref{eq:update} ;
  \item une cellule \emph{extérieure} ($V_i\cap D=\emptyset$) disparaît du tissu ; son germe est \emph{figé}. Les
    germes figés forment la paroi qui contient le tissu : sans eux, les cellules du bord s'étaleraient
    jusqu'au rectangle ;
  \item une cellule \emph{coupée} va vers le centroïde $\vect c_i$ de sa partie intérieure, avec un poids
    proportionnel à son aire extérieure $A^{\mathrm{out}}_i=|V_i|-|V_i\cap D|$ :
    \begin{equation}\label{eq:beta}
      \vect p_i'=\frac{\vect c_i+\alpha\,\vect b+w_i\,\vect c_i}{1+\alpha+w_i},\qquad w_i=\beta\,A^{\mathrm{out}}_i .
    \end{equation}
    Avec $\beta$ grand (défaut $10^6$), $\vect p_i'\simeq\vect c_i$ : les cellules du bord suivent Lloyd sans être
    attirées, ce qui garde un bord régulier ; avec $\beta=0$ elles sont attirées comme les autres
    (figure~\ref{fig:domaines}). Pour une cellule intérieure, $w_i=0$ et \eqref{eq:beta} se réduit
    exactement à~\eqref{eq:update}.
\end{itemize}
Le rectangle doit déborder un peu de la forme (une ou deux cellules) pour qu'il reste des germes extérieurs.
C'est la méthode de génération d'échantillons développée en Python avec \code{shapely} ; \cy{} la reprend sans
dépendance, avec une forme quelconque au lieu d'un rectangle.

\textbf{Découpe d'une cellule} (figure~\ref{fig:clip}). Le contour de la cellule convexe est échantillonné au pas
$h=s/4$ et $f$ est évaluée en chaque échantillon. Si tous sont intérieurs, la cellule est intérieure ; si tous sont
extérieurs, elle est extérieure. Sinon, en partant d'un échantillon intérieur, on parcourt le contour : les sommets
intérieurs sont gardés ; chaque changement de signe de $f$ entre deux échantillons donne un point de passage,
calculé par dichotomie (60 itérations, précision machine) ; entre une \emph{sortie} et l'\emph{entrée} suivante,
le bord de la forme traverse la cellule : des points de la corde, espacés de $h/2$, sont projetés sur $f=0$ par
la méthode de Newton ($\vect x\leftarrow\vect x-f\,\nabla f/\|\nabla f\|^2$, gradient par différences centrées) et
ajoutés s'ils restent dans la cellule. Le bord courbe est ainsi suivi finement, et la partie intérieure peut être
non convexe (bord concave). Un point de passage situé sur une paroi commune à deux cellules est calculé de la même
façon des deux côtés : il devient un sommet partagé.

\begin{figure}[H]
  \centering
  \begin{tikzpicture}[scale=1.25]
    \def\hex{(2,0)--(1,1.732)--(-1,1.732)--(-2,0)--(-1,-1.732)--(1,-1.732)--cycle}
    \begin{scope}
      \clip \hex;
      \fill[accent!12] (-4.2,0.4) circle (4.9);
    \end{scope}
    \draw[thick] \hex;
    \draw[bundle,very thick] (-4.2,0.4) ++(-28:4.9) arc (-28:28:4.9) node[above right,font=\footnotesize]{$f=0$};
    \node[font=\footnotesize] at (-0.8,0.2) {$f<0$ (tissu)};
    \node[font=\footnotesize] at (2.6,1.0) {$f>0$};
    \foreach \p in {{(0.500,1.732)},{(0.000,1.732)},{(-0.500,1.732)},{(-1.000,1.732)},{(-1.250,1.299)},{(-1.500,0.866)},
        {(-1.750,0.433)},{(-2.000,0.000)},{(-1.750,-0.433)},{(-1.500,-0.866)},{(-1.250,-1.299)},{(-1.000,-1.732)},
        {(-0.500,-1.732)},{(-0.000,-1.732)}} \fill \p circle (1.4pt);
    \foreach \p in {{(2.000,0.000)},{(1.750,0.433)},{(1.500,0.866)},{(1.250,1.299)},{(1.000,1.732)},{(0.500,-1.732)},
        {(1.000,-1.732)},{(1.250,-1.299)},{(1.500,-0.866)},{(1.750,-0.433)}} \draw[fill=white] \p circle (1.4pt);
    \draw[gray,dashed] (0.515,1.732)--(0.212,-1.732);
    \foreach \a/\b in {{(0.494,1.485)}/{(0.574,1.503)},{(0.450,0.990)}/{(0.661,1.016)},{(0.407,0.495)}/{(0.699,0.501)},
        {(0.364,0.000)}/{(0.681,-0.028)},{(0.320,-0.495)}/{(0.607,-0.552)},{(0.277,-0.990)}/{(0.480,-1.053)},
        {(0.234,-1.485)}/{(0.309,-1.517)}} {\draw[->,gray,>=Stealth] \a--\b; \fill[accent] \b circle (1.3pt);}
    \fill[bundle] (0.515,1.732) circle (2pt) node[above left,font=\footnotesize,text=black]{sortie};
    \fill[bundle] (0.212,-1.732) circle (2pt) node[below left,font=\footnotesize,text=black]{entrée};
    \begin{scope}[shift={(4.3,0.9)},font=\footnotesize]
      \fill (0,0.6) circle (1.4pt); \node[right] at (0.15,0.6) {échantillon intérieur};
      \draw[fill=white] (0,0.2) circle (1.4pt); \node[right] at (0.15,0.2) {échantillon extérieur};
      \fill[bundle] (0,-0.2) circle (2pt); \node[right] at (0.15,-0.2) {passage (dichotomie)};
      \fill[accent] (0,-0.6) circle (1.4pt); \node[right] at (0.15,-0.6) {projection de Newton};
      \fill[accent!12] (-0.1,-1.1) rectangle (0.1,-0.9); \node[right] at (0.15,-1.0) {partie intérieure};
    \end{scope}
  \end{tikzpicture}
  \caption{Découpe d'une cellule de Voronoï convexe par le bord de la forme (schéma ; une projection sur deux est
    représentée, et le pas d'échantillonnage est exagéré).}
  \label{fig:clip}
\end{figure}

\textbf{Écrire une forme.} \code{shape} est une expression de $f$ : opérateurs \code{+ - * / \textasciicircum}
(puissance), parenthèses, fonctions \code{sin cos tan asin acos atan exp log sqrt abs floor ceil} et
\code{atan2(y,x) min(a,b) max(a,b) pow(a,b)}, constantes \code{pi} et \code{e}. Variables : \code{x}, \code{y} ;
\code{r} et \code{t}, rayon et angle (radians, dans $]-\pi,\pi]$) depuis le centre du rectangle ; \code{cx},
\code{cy} (centre), \code{Lx}, \code{Ly}. Une fonction distance signée (comme \code{r - 18}) est idéale, mais
n'importe quelle fonction continue convient pourvu que son gradient ne s'annule pas sur le bord. Opérations
booléennes : union \code{min(f1, f2)}, intersection \code{max(f1, f2)}, différence \code{max(f1, -f2)}.

\begin{center}
\begin{tabular}{@{}ll@{}}
  \toprule
  Forme & \code{shape} \\
  \midrule
  disque de rayon 18 & \code{r - 18} \\
  ellipse de demi-axes 18 et 11 & \code{((x-cx)/18)\textasciicircum2 + ((y-cy)/11)\textasciicircum2 - 1} \\
  rectangle centré $30\times16$ & \code{max(abs(x-cx) - 15, abs(y-cy) - 8)} \\
  anneau & \code{max(r - 18, 9 - r)} \\
  disque échancré à droite & \code{r - 20 + 4.5*exp(-t\textasciicircum2/0.12)} \\
  deux disques réunis & \code{min(sqrt((x-14)\textasciicircum2+(y-cy)\textasciicircum2) - 10, sqrt((x-30)\textasciicircum2+(y-cy)\textasciicircum2) - 10)} \\
  \bottomrule
\end{tabular}
\end{center}

\begin{figure}[H]
  \centering
  \begin{subfigure}{0.32\linewidth}\includegraphics[width=\linewidth]{figures/dom_disque.png}\caption{\code{r - 18}}\end{subfigure}\hfill
  \begin{subfigure}{0.32\linewidth}\includegraphics[width=\linewidth]{figures/dom_anneau.png}\caption{anneau}\end{subfigure}\hfill
  \begin{subfigure}{0.32\linewidth}\includegraphics[width=\linewidth]{figures/dom_ellipse.png}\caption{ellipse, $K=1000$, angle $0^\circ$}\end{subfigure}\\[2mm]
  \begin{subfigure}{0.4\linewidth}\includegraphics[width=\linewidth]{figures/rein.png}\caption{\code{examples/rein.txt}}\end{subfigure}\hfill
  \begin{subfigure}{0.26\linewidth}\includegraphics[width=\linewidth]{figures/rein_zoom.png}\caption{zoom sur l'échancrure}\end{subfigure}\hfill
  \begin{subfigure}{0.3\linewidth}\includegraphics[width=\linewidth]{figures/rein_beta0.png}\caption{même forme, $\beta=0$}\end{subfigure}
  \caption{Domaines définis par une fonction ($\alpha=0.05$, 5\,\% de bundles, 50 itérations). Le bord du tissu est
    tracé en trait épais.}
  \label{fig:domaines}
\end{figure}

\textbf{Limites.} Une partie de la forme plus petite qu'une cellule et qui ne coupe aucun de ses côtés
(îlot isolé) n'est pas vue. Un bord très sinueux à l'échelle d'une cellule peut couper un côté plusieurs fois
entre deux échantillons ($h=s/4$) ; la forme doit varier lentement à cette échelle. Une cellule coupée en deux
morceaux disjoints est représentée par un seul polygone.

\subsection{Pavage à sommets partagés}\label{sec:mesh}

Chaque cellule est calculée indépendamment : un sommet commun à trois cellules est donc calculé trois fois, à
quelques $10^{-16}$ près. Les sommets sont fusionnés à la tolérance $10^{-8}s$ (table de hachage sur une grille de
pas $10^{-6}s$). Une paroi qui n'appartient qu'à une cellule est une paroi du \emph{bord du tissu} ; une cellule qui
en possède une est une cellule de bord. Ce critère est purement topologique : il vaut pour le rectangle comme pour
n'importe quelle forme.

\subsection{Métriques}\label{sec:metriques}

Calculées sur le pavage final (cellules du tissu, bords compris), elles reprennent celles de l'article :
\begin{center}
\begin{tabularx}{\linewidth}{@{}lX@{}}
  \toprule
  Nom & Définition \\
  \midrule
  \code{nCells}, \code{nBundles} & nombre de cellules du tissu, de bundles \\
  \code{meanSI} & indice de forme moyen $\langle SI\rangle=\langle4\pi A/P^2\rangle$ (compacité ; 1 pour un disque,
    0.907 pour l'hexagone régulier, 0.785 pour le carré) \\
  \code{cvA} & coefficient de variation des aires $\sigma_A/\langle A\rangle$ (hétérogénéité) \\
  \code{dNN} & distance moyenne entre le centroïde d'une cellule et le plus proche autre centroïde, rapportée à $s$
    (environ 1.07 pour l'hexagone régulier) \\
  \code{dM} & distance de Mahalanobis moyenne entre le centroïde d'une cellule ordinaire et le plus proche centroïde de
    bundle, rapportée à $s$ (avec $\Sigma$ telle que définie, non normalisée) \\
  \code{areaRegular}, \code{areaBundle} & aire moyenne des cellules ordinaires, des bundles ($A_0=0.785$) \\
  \code{SIRegular}, \code{SIBundle} & indice de forme moyen des cellules ordinaires, des bundles \\
  \code{meanSides} & nombre moyen de côtés des cellules intérieures (6 en moyenne pour un pavage de Voronoï,
    par le théorème d'Euler) \\
  \code{lastMove} & déplacement moyen des germes à la dernière itération, rapporté à $s$ \\
  \bottomrule
\end{tabularx}
\end{center}
La ligne de commande affiche aussi l'histogramme du nombre de côtés des cellules intérieures.

%======================================================================
\section{Fichiers}\label{sec:fichiers}
%======================================================================

\subsection{Fichier de cellules}

Fichier texte (\code{\#} pour les commentaires), écrit avec 17 chiffres significatifs :
\begin{lstlisting}
cellyod 1
param <mot-clé> <valeur>        tous les paramètres du tissu (même syntaxe que le fichier de paramètres)
domain xmin ymin xmax ymax       rectangle des germes
iteration <n>                   nombre d'itérations effectuées
vertices <N>                    puis N lignes : x y
cells <M>                       puis M lignes : type sx sy nv v0 v1 ... v(nv-1)
frozen <K>                      puis K lignes : sx sy   (germes extérieurs ; absent sans forme)
\end{lstlisting}
L'identifiant d'une cellule est son rang dans la section \code{cells} (à partir de 0) ; \code{type} vaut 0 pour une
cellule ordinaire, 1 pour un bundle ; \code{(sx, sy)} est son germe ; ses \code{nv} sommets, indices dans la section
\code{vertices}, sont donnés dans le sens trigonométrique. Les sommets sont partagés : deux cellules voisines
référencent les mêmes indices le long de leur paroi commune, et une paroi n'appartenant qu'à une cellule est sur
le bord du tissu. Exemple (réseau hexagonal $6\times6$, début du fichier) :
\lstinputlisting{figures/petit_extrait.txt}

\paragraph{Lire le fichier.} Pour un autre programme : lire les lignes \code{param} si besoin, puis $N$ couples
$(x,y)$ après \code{vertices}, puis $M$ lignes après \code{cells} dont on garde \code{type} et la liste de sommets.
Une représentation « un polygone par cellule » s'en déduit directement (une option d'export dans ce sens est
prévue).

\subsection{Reprise d'un calcul}\label{sec:reprise}

Le pavage est entièrement déterminé par les germes et les paramètres : \code{-load} (ou \emph{Charger} dans
l'interface) relit les paramètres, les germes des cellules et les germes figés, et reconstruit le pavage. Comme le
calcul de Voronoï ne dépend pas de la numérotation des germes, poursuivre un calcul enregistré donne exactement
(au bit près) les mêmes germes qu'un calcul d'une traite. Seule la numérotation des cellules peut différer quand
une forme est utilisée : un germe figé dont la cellule entre dans la forme au cours du calcul est placé à la fin.

\subsection{Fichier \texttt{morphology}}\label{sec:morpho}

Avec \code{morphology fichier}, chaque exécution ajoute une ligne : tous les paramètres numériques et de choix
(la forme, texte, n'y figure pas), le nombre d'itérations, puis les métriques de la section~\ref{sec:metriques}.
L'en-tête (\code{\#} puis les noms de colonnes) est écrit à la création du fichier. Un balayage se fait en
relançant \cy{} avec des paramètres différents et le même fichier (section~\ref{sec:etudes}).

\subsection{SVG}

Vue vectorielle du tissu, dans le même style que l'interface : parois fines, bord du tissu épais, germes gris
cerclés, bundles en rouge, légende. Deux thèmes : \code{svgTheme light} (défaut, noir sur blanc : cellules blanches,
parois noires) et \code{svgTheme dark} (blanc sur noir : cellules sombres sur fond noir, parois blanches) ; les
couleurs des deux thèmes sont définies dans \code{src/theme.hpp}, commun au SVG et à l'interface.
\code{svgSeeds 0} retire les germes,
\code{svgLegend 0} la légende ; \code{svg none} supprime le fichier. Un SVG s'ouvre dans un navigateur (zoom
illimité) et se convertit en PNG avec Inkscape :
\code{inkscape cells.svg --export-type=png -w 2000}.

%======================================================================
\section{Interface graphique}\label{sec:gui}
%======================================================================

\begin{figure}[H]
  \centering
  \includegraphics[width=\linewidth]{figures/gui.png}
  \caption{\cyg{} avec \code{examples/rein.txt}. À gauche la vue, à droite le panneau (qui défile), en bas la
    barre d'état.}
  \label{fig:gui}
\end{figure}

\textbf{Vue.} Molette : zoom autour de la souris ; glisser (n'importe quel bouton) : déplacer ; clic gauche sans
déplacement sur une cellule : bundle on/off. La barre d'état donne les coordonnées et la cellule survolée
(numéro, nombre de côtés, aire, indice de forme), l'aide du paramètre survolé dans le panneau, ou le dernier
message (erreurs en rouge).

\textbf{Panneau.}
\begin{itemize}
  \item \emph{Simulation} : \emph{Lancer/Pause} (Espace), \emph{Pas} (S ou $\rightarrow$, une itération),
    \emph{Régénérer} (R, retour à l'itération 0 avec de nouveaux germes), \emph{Calcul direct} (sans animation,
    jusqu'à \code{lloydIterations}) ; arrêt automatique à \code{lloydIterations} (à décocher pour itérer sans
    limite) ; nombre d'itérations par image (l'animation reste limitée à environ 30~ms de calcul par image).
  \item \emph{Paramètres} : tous les paramètres, groupés comme dans \code{-h}. Un curseur (logarithmique pour
    \code{anisotropy} et \code{beta}) et une valeur ; cliquer sur la valeur pour la taper (Entrée pour valider).
    Les choix (\code{lattice}, \code{bundleMotion}) sont des boutons ; \code{shape} est un champ de texte. Un
    paramètre marqué \code{*} régénère les germes et relance l'animation ; les autres s'appliquent à l'itération
    suivante (section~\ref{sec:effets}).
  \item \emph{Affichage} : \emph{Uni} (rendu de la figure~\ref{fig:gui}), ou coloration par nombre de côtés
    (palette discrète ; les cellules de bord sont estompées), par aire ou par indice de forme (palette continue,
    échelle min--max affichée) ; germes et parois ; thème \emph{Noir sur blanc} (défaut) ou \emph{Blanc sur noir} (touche T), appliqué
    aussi au SVG écrit par \emph{Enregistrer} ; \emph{Recadrer} (F).
  \item \emph{Métriques} : valeurs courantes ; un clic sur une métrique trace son évolution en fonction de
    l'itération depuis la dernière régénération ; histogramme du nombre de côtés.
  \item \emph{Fichiers} : nom du fichier de cellules ; \emph{Enregistrer} (Ctrl+S) écrit ce fichier et le SVG de
    même nom ; \emph{Charger} relit un fichier de cellules ou de paramètres ; \emph{Paramètres} écrit
    \code{<nom>\_params.txt}. Un fichier peut aussi être glissé-déposé dans la fenêtre.
\end{itemize}
Les réglages d'affichage (taille de fenêtre, police, coloration, thème, vitesse\dots) sont conservés dans
\code{cellyod-gui.ini}, dans le répertoire courant. \code{cellyod-gui fichier --screenshot image.png} calcule
jusqu'à \code{lloydIterations}, enregistre une capture et quitte (utilisé pour la figure~\ref{fig:gui}).

\textbf{Démarche conseillée.} (1) Régler le domaine et la forme avec $\alpha=0$ et quelques itérations ;
(2) choisir \code{bundleFraction} ; (3) lancer l'animation et ajuster $\alpha$, $K$ et l'angle en regardant le tissu
et la courbe de \code{cvA} ; (4) arrêter au bon nombre d'itérations, ajouter ou retirer des bundles à la souris
si besoin ; (5) enregistrer, et exporter les paramètres pour reproduire ou décliner le tissu en ligne de
commande.

%======================================================================
\section{Études types}\label{sec:etudes}
%======================================================================

\paragraph{Balayage d'un paramètre.} Les courbes de la figure~\ref{fig:curves} sont produites ainsi
(\code{doc/make\_figures.sh}) :
\begin{lstlisting}
for a in 0 0.01 0.02 0.05 0.1 0.15 0.2; do
  ./cellyod examples/tissu.txt alpha=$a output=/dev/null svg=none morphology=balayage.txt
done
\end{lstlisting}
Pour une évolution en fonction des itérations, varier \code{lloydIterations} de la même façon, ou enchaîner des
reprises (\code{-load}) pour ne pas recalculer le début.

\paragraph{Variabilité.} Les métriques dépendent du tirage : répéter avec plusieurs \code{seed} et moyenner.
Pour le tissu de référence en $30\times30$ et 50 itérations, six graines donnent \code{cvA} entre 0.63 et 0.75
(écart type d'environ 0.04), mais $\langle SI\rangle$ entre 0.862 et 0.865 : l'hétérogénéité des aires, dominée par
quelques amas, est la métrique la plus sensible au tirage, et un domaine plus grand la réduit.

\paragraph{Régimes successifs.} Les paramètres vivants permettent d'enchaîner des régimes : par exemple
20 itérations de Lloyd pur pour régulariser, puis l'attraction :
\begin{lstlisting}
./cellyod lattice=uniform lloydIterations=20 output=a.txt svg=none
./cellyod -load a.txt bundleFraction=0.05 alpha=0.05 lloydIterations=70   # ! bundleFraction ignoré
\end{lstlisting}
Attention : un paramètre de régénération (ici \code{bundleFraction}) n'a pas d'effet sur une reprise ; dans ce cas,
choisir les bundles dès le départ avec $\alpha=0$, puis passer $\alpha$ à 0.05 à la reprise.

\paragraph{Grands tissus.} Le temps est proportionnel au nombre de cellules (environ 2~µs par cellule et par
itération, plus la recherche du bundle le plus proche, proportionnelle au nombre de bundles) : $100\times100$
(12\,700 cellules), 20 itérations, en 0.7~s.

%======================================================================
\section{Organisation du code}\label{sec:code}
%======================================================================

\begin{center}
\begin{tabularx}{\linewidth}{@{}lX@{}}
  \toprule
  Fichier & Rôle \\
  \midrule
  \code{src/geom.hpp} & vecteur \code{V}, polygones : aire signée, périmètre, centroïde \\
  \code{src/theme.hpp} & thèmes de rendu (noir sur blanc, blanc sur noir), communs au SVG et à l'interface \\
  \code{src/expr.hpp/.cpp} & analyseur d'expressions (descente récursive vers une notation polonaise inverse,
    évaluée sur une pile) \\
  \code{src/params.hpp/.cpp} & structure \code{Params} et registre \code{paramSpecs()} ; lecture et écriture des
    paramètres, aide \\
  \code{src/tissue.hpp/.cpp} & \code{Domain} (rectangle, forme, découpe), \code{Mahalanobis}, \code{Tissue}
    (germes, Voronoï, itérations, sommets partagés), \code{computeMetrics} \\
  \code{src/io.hpp/.cpp} & fichier de cellules, SVG, fichier \code{morphology} \\
  \code{src/cli.cpp} & ligne de commande \\
  \code{gui/cellyod\_gui.cpp} & interface raylib/raygui \\
  \code{doc/} & ce tutoriel et \code{make\_figures.sh} \\
  \bottomrule
\end{tabularx}
\end{center}

\textbf{Classe \code{Tissue}.} \code{init(P)} place les germes et choisit les bundles ; \code{step()} fait une
itération (avec les paramètres courants \code{P}, que l'on peut modifier entre deux pas) ; \code{run()} itère jusqu'à
\code{lloydIterations} ; \code{restore()} reprend un état enregistré ; \code{mesh()} construit le pavage à sommets
partagés. Les membres \code{seeds}, \code{isBundle}, \code{voronoi} (cellules dans le rectangle) et \code{cells}
(parties dans la forme, vides pour les germes figés) décrivent l'état courant.

\textbf{Ajouter un paramètre.} Un champ dans \code{Params} et une ligne dans \code{paramSpecs()} (nom, groupe, aide,
bornes du curseur, échelle logarithmique, régénération ou non) suffisent : la lecture des fichiers, l'écriture dans
les fichiers de cellules, l'aide \code{-h}, le fichier \code{morphology} et le panneau de l'interface sont construits
à partir du registre. Il reste à l'utiliser dans \code{Tissue::init} ou \code{Tissue::step}.

\textbf{Ajouter une métrique.} Un champ dans \code{Metrics}, son calcul dans \code{computeMetrics}, une entrée
dans \code{Metrics::items()} : elle apparaît dans l'interface (avec sa courbe) et dans le fichier \code{morphology}.

\textbf{Ajouter une coloration.} Une ligne dans \code{COLOR\_FIELDS} (\code{gui/cellyod\_gui.cpp}) : nom,
discrète ou continue, fonction donnant la valeur d'une cellule.

\textbf{Changer la forme du domaine.} Tout passe par \code{Domain} : \code{f} (valeur de la fonction),
\code{inShape} et \code{clip} (partie d'une cellule dans la forme). Une autre description de la forme (polygone lu
dans un fichier, image binaire\dots) se branche en remplaçant l'évaluation de $f$.

\textbf{Validation.} Sans forme, \cy{} reproduit \code{nodeGenPlantCom} (avec \code{minEdgeRatio 0}, c'est-à-dire sans
reconformation) : mêmes nombres de cellules et mêmes métriques pour les quatre états initiaux, avec ou sans
bundles et anisotropie. Les reprises (\code{-load}) sont vérifiées identiques au bit près.

\textbf{Pistes.} Export « un polygone par cellule » ; décalage des parois (épaisseur) pour les codes
d'éléments discrets ; forme lue dans un fichier (polygone, image) ; métriques par région ; parallélisation du
calcul de Voronoï (chaque cellule est indépendante).

\end{document}
